\section{Summary of the Default ARDs Main Characteristics}
\label{sec:prod_summ}

The generation of the proposed products can follow different processing flows, depending on the availability of the input data described in the previous sections, the mission (e.g., Sentinel-1 or ROSE-L), the desired timeliness, the application domain, and other factors.

In this section, the default characteristics of each product are provided, including sampling, applied corrections, input data type, and other product-specific properties.

\subsection{CSLC}
\label{sec:cslc_summ}

\begin{longtblr}[
  label = {tab:sum_cslc},
  entry = none,
  caption = {Main characteristics of the default \acs{cslc} product.},
]{
  width   = \linewidth,
  colspec = {X[l] X[l] X[c] X[c]},
  colsep  = 2pt,
  hlines,
  hline{1-Z} = {-}{0.05em},
}

\SetCell[r=2]{c}\textbf{Characteristic Item}
  & \SetCell[r=2]{c}\textbf{Input Type}
  & \SetCell[c=2]{c}\textbf{Mission} & \\
& & \textbf{Sentinel-1} & \textbf{ROSE-L} \\

Radar Reference Grid & & Virtual & Virtual \\

Geometric alignment & { + COP\acs{dem} 30 m \\
                        + Precise Orbits } & \ding{51} & \ding{51} \\

Topographic phase correction & COP\acs{dem} 30 m & \ding{51} & \ding{51} \\

Atmospheric delay \& timing corrections & \acs{etad} & \ding{51} & \ding{51} \\

ESD & Geometric-corrected \acs{cslc}s (GC-CSLC) & - & \ding{51} \\

ICC & Azimuth-shift-corrected GC-\acs{cslc}s (ASGC-CSLC) & - & \ding{51} \\

Split-Spectrum & { + Range-Azimuth-shift-corrected GC-\acs{cslc}s (RASGC-CSLC) \\
                   + Reference ASGC-\acs{cslc} }& - & \ding{51} \\

\SetCell[r=2]{l}Spacing
  & \SetCell[r=2]{l}
  & \SetCell[r=2,c=2]{c}
    {Defined by the Radar Reference Grid:\\
     minimum between Sentinel-1\\
     and ROSE-L spacing} & \\
& & & \\

Multilooking &  & - &  - \\

Denoising &  & - &  - \\

Slicing &  & S-1 burst &  ROSE-L burst \\

Reference System &  & Azimuth/Slant-range &  Azimuth/Slant-range \\

\end{longtblr}

Note that the inputs for \acs{esd}, \acs{icc}, and split-spectrum all originate from a progressively refined form of the \acs{cslc} product. Specifically, ESD applies azimuth corrections starting from the \acs{cslc} obtained through geometric coregistration (referred to as GC-CSLC in the previous table). \acs{icc} then applies range corrections to the output of \acs{esd} (referred to as ASGC-CSLC), and finally, the split-spectrum technique provides the ultimate refinement, applied to the output of \acs{icc} (referred to as RASGC-CSLC in the previous table).

\subsection{GSLC}
\label{sec:gslc_summ}

The default \acs{gslc} product is generated from the \acs{cslc}, from which it inherits the main characteristics defined in \ref{sec:cslc_summ}. Below, only the product-specific main characteristics are reported.

\begin{longtblr}[
  label = {tab:sum_gslc},
  entry = none,
  caption = {Main characteristics of the default \acs{gslc} product.},
]{
  width   = \linewidth,
  colspec = {X[l] X[l] X[c] X[c]},
  colsep  = 2pt,
  hlines,
  hline{1-Z} = {-}{0.05em},
}

\SetCell[r=2]{c}\textbf{Characteristic Item}
  & \SetCell[r=2]{c}\textbf{Input Type}
  & \SetCell[c=2]{c}\textbf{Mission} & \\
& & \textbf{Sentinel-1} & \textbf{ROSE-L} \\

Amplitude convention & \acs{cslc} & beta0 & beta0 \\

Phase & Point-based interpolation from \acs{cslc} phase & \ding{51} & \ding{51} \\

Spacing & & {Comparable with the \\
            \acs{cslc} ground spacing} & 
            {Comparable with the \\
            \acs{cslc} ground spacing} \\

Multilooking &  & - &  - \\

Slicing &  & S-1 projected burst &  ROSE-L projected burst \\

Reference System &  & Cartographic: UTM/UPS &  Cartographic: UTM/UPS \\

\end{longtblr}

\subsection{NRB}
\label{sec:nrb_summ}

The default \acs{nrb} product is generated from the \acs{cslc}. No timing or atmospheric corrections are applied during \acs{nrb} generation, as these are handled at the \acs{cslc} level. Below, only the product-specific main characteristics are reported.

\begin{longtblr}[
  label = {tab:sum_nrb},
  entry = none,
  caption = {Main characteristics of the default \acs{nrb} product.},
]{
  width   = \linewidth,
  colspec = {X[l] X[l] X[c] X[c]},
  colsep  = 2pt,
  hlines,
  hline{1-Z} = {-}{0.05em},
}

\SetCell[r=2]{c}\textbf{Characteristic Item}
  & \SetCell[r=2]{c}\textbf{Input Type}
  & \SetCell[c=2]{c}\textbf{Mission} & \\
& & \textbf{Sentinel-1} & \textbf{ROSE-L} \\

Amplitude convention &  & gamma0 & gamma0 \\

\acs{rtc} & { + \acs{cslc} \\
              + COP\acs{dem} 30} m & \ding{51} & \ding{51} \\

Shiroma for \acs{rtc} & & \ding{51} & \ding{51} \\

\acs{anf} & & \ding{51} & \ding{51} \\

Local Incidence Angle & & \ding{51} & \ding{51} \\

Spacing & & 10 m x 10 m & 10 m x 10 m \\

Multilooking &  & - &  - \\

Slicing &  & US-\acs{mgrs} 100 km x 100 km &  US-\acs{mgrs} 100 km x 100 km \\

Reference System &  & Cartographic: UTM/UPS &  Cartographic: UTM/UPS \\

\end{longtblr}

\subsection{POL}
\label{sec:pol_summ}

The default \acs{pol} product is generated from the \acs{cslc}. No timing or atmospheric corrections are applied during \acs{pol} generation, as these are handled at the \acs{cslc} level. Below, only the product-specific main characteristics are reported.

Multi-looking is carried out for a proper covariance matrix estimation, but the number of looks (and hence the estimation window size) is defined by the user. A typical multi-look size is provided in \autoref{tab:sum_pol}. As the multi-looking degrades the spatial resolution, the output pixel spacing can be enlarged, e.g. to 20 m x 20 m or 30 m x 30 m, to reduce the file size of the product.

\begin{longtblr}[
  label = {tab:sum_pol},
  entry = none,
  caption = {Main characteristics of the default \acs{pol} product.},
]{
  width   = \linewidth,
  colspec = {X[l] X[l] X[c] X[c]},
  colsep  = 2pt,
  hlines,
  hline{1-Z} = {-}{0.05em},
}

\SetCell[r=2]{c}\textbf{Characteristic Item}
  & \SetCell[r=2]{c}\textbf{Input Type}
  & \SetCell[c=2]{c}\textbf{Mission} & \\
& & \textbf{Sentinel-1} & \textbf{ROSE-L} \\

Amplitude convention &  & gamma0 & gamma0 \\

\acs{rtc} & { + \acs{cslc} \\
              + COP\acs{dem} 30} m & \ding{51} & \ding{51} \\

Shiroma for \acs{rtc} & & \ding{51} & \ding{51} \\

\acs{anf} & & \ding{51} & \ding{51} \\

Local Incidence Angle & & \ding{51} & \ding{51} \\

Spacing & & 10 m x 10 m & 10 m x 10 m \\

Multilooking (az x rg) &  & 3 x 15 & 5 x 10 \\

Slicing &  & US-\acs{mgrs} 100 km x 100 km &  US-\acs{mgrs} 100 km x 100 km \\

Reference System &  & Cartographic: UTM/UPS &  Cartographic: UTM/UPS \\

\end{longtblr}


\subsection{ORB}
\label{sec:orb_summ}

The default \acs{orb} product is generated from the \acs{cslc}. No timing or atmospheric corrections are applied during \acs{orb} generation, as these are handled at the \acs{cslc} level. Below, only the product-specific main characteristics are reported.

A typical multi-look size is provided in \autoref{tab:sum_orb}. As the multi-looking degrades the spatial resolution, the output pixel spacing can be enlarged, e.g. to 300 m x 300 m, to reduce the file size of the product.

\begin{longtblr}[
  label = {tab:sum_orb},
  entry = none,
  caption = {Main characteristics of the default \acs{orb} product.},
]{
  width   = \linewidth,
  colspec = {X[l] X[l] X[c] X[c]},
  colsep  = 2pt,
  hlines,
  hline{1-Z} = {-}{0.05em},
}

\SetCell[r=2]{c}\textbf{Characteristic Item}
  & \SetCell[r=2]{c}\textbf{Input Type}
  & \SetCell[c=2]{c}\textbf{Mission} & \\
& & \textbf{Sentinel-1} & \textbf{ROSE-L} \\
Intensity convention & CSLC & sigma0 & sigma0 \\
\acs{rtc} & { + \acs{cslc} \\
              + COP\acs{dem} 30} m & \ding{51} & \ding{51} \\

Shiroma for \acs{rtc} & & \ding{51} & \ding{51} \\

\acs{anf} & & \ding{51} & \ding{51} \\
Local Incidence Angle & & \ding{51} & \ding{51} \\
Spacing & & 10 m x 10 m & 10 m x 10 m \\

Multilooking (az x rg) &  & 300 m x 300 m & 300 m x 300 m \\

Slicing &  & US-\acs{mgrs} 100 km x 100 km &  US-\acs{mgrs} 100 km x 100 km \\

Reference System &  & Cartographic: UTM/UPS &  Cartographic: UTM/UPS \\

\end{longtblr}

\subsection{PRO/PAO}
\label{sec:pro_pao_summ}

The default \acs{pro} \& \acs{pao} products are generated from the \acs{cslc}. No timing or atmospheric corrections are applied during \acs{pro} \& \acs{pao} generation, as these are handled at the \acs{cslc} level. Below, only the product-specific main characteristics are reported.

\begin{longtblr}[
  label = {tab:sum_pro_pao},
  entry = none,
  caption = {Main characteristics of the default \acs{pro} \& \acs{pao} products.},
]{
  width   = \linewidth,
  colspec = {X[l] X[l] X[c] X[c]},
  colsep  = 2pt,
  hlines,
  hline{1-Z} = {-}{0.05em},
}

\SetCell[r=2]{c}\textbf{Characteristic Item}
  & \SetCell[r=2]{c}\textbf{Input Type}
  & \SetCell[c=2]{c}\textbf{Mission} & \\
& & \textbf{Sentinel-1} & \textbf{ROSE-L} \\

Estimated Offsets & Pairs of \acs{cslc} & Along slant-range and azimuth &  Along slant-range and azimuth \\

Pixel offset algorithm  &  & Amplitude CC &  Amplitude CC \\

Spacing & & ~50 m x ~50 m & ~50 m x ~50 m \\

Slicing &  & S-1 burst &  ROSE-L burst \\

Reference System &  & Azimuth/Slant-range &  Azimuth/Slant-range \\

\end{longtblr}

\subsection{GPRO/GPAO}
\label{sec:gpro_gpao_summ}

The default \acs{gpro} \& \acs{gpao} are the geocoded \acs{gpro} \& \acs{gpao} products, from which they inherit the main characteristics defined in \ref{sec:pro_pao_summ}. Below, only the product-specific main characteristics are reported.

\begin{longtblr}[
  label = {tab:sum_gpro_gpao},
  entry = none,
  caption = {Main characteristics of the default \acs{gpro} \& \acs{gpao} products.},
]{
  width   = \linewidth,
  colspec = {X[l] X[l] X[c] X[c]},
  colsep  = 2pt,
  hlines,
  hline{1-Z} = {-}{0.05em},
}

\SetCell[r=2]{c}\textbf{Characteristic Item}
  & \SetCell[r=2]{c}\textbf{Input Type}
  & \SetCell[c=2]{c}\textbf{Mission} & \\
& & \textbf{Sentinel-1} & \textbf{ROSE-L} \\

Estimated Offsets & Pairs of \acs{cslc} & Along slant-range and azimuth &  Along slant-range and azimuth \\

Ground Pojection  & { + COP\acs{dem} 30m \\
                      + Precise Reference Orbit}& \ding{51} & \ding{51} \\

Spacing & & 50 m x 50 m & 50 m x 50 m \\

Slicing &  & S-1 projected burst &  ROSE-L projected burst \\

Reference System &  & Cartographic: UTM/UPS &  Cartographic: UTM/UPS \\

\end{longtblr}

\subsection{INT/COHE/UNW}
\label{sec:int_summ}

The default \acs{int}, \acs{cohe} and \acs{unw} products are generated from the \acs{cslc}. No timing or atmospheric corrections are applied during the  \acs{int}, \acs{cohe} and \acs{unw} product generation, as these are handled at the \acs{cslc} level. Below, only the product-specific main characteristics are reported.

\begin{longtblr}[
  label = {tab:sum_int},
  entry = none,
  caption = {Main characteristics of the default \acs{int}, \acs{cohe} and \acs{unw} products.},
]{
  width   = \linewidth,
  colspec = {X[l] X[l] X[c] X[c]},
  colsep  = 2pt,
  hlines,
  hline{1-Z} = {-}{0.05em},
}

\SetCell[r=2]{c}\textbf{Characteristic Item}
  & \SetCell[r=2]{c}\textbf{Input Type}
  & \SetCell[c=2]{c}\textbf{Mission} & \\
& & \textbf{Sentinel-1} & \textbf{ROSE-L} \\

Wrapped interferogram & Pairs of \acs{cslc} & \ding{51} &  \ding{51} \\

Interferometric coherence & Pairs of \acs{cslc} & \ding{51} &  \ding{51} \\

Unwrapped interferogram & { + Wrapped interferogram \\
                            + Interferometric coherence } & \ding{51} &  \ding{51} \\

Spacing & & ~50 m x ~50 m & ~50 m x ~50 m \\

Slicing &  & S-1 burst &  ROSE-L burst \\

Reference System &  & Azimuth/Slant-range &  Azimuth/Slant-range \\

\end{longtblr}

\subsection{GINT/GCOHE/GUNW}
\label{sec:gint_summ}

The default \acs{gint}, \acs{gcohe} and \acs{gunw} are the geocoded \acs{int}, \acs{cohe} and \acs{unw} products, respectively, from which they inherit the main characteristics defined in \ref{sec:int_summ}. Below, only the product-specific main characteristics are reported.

\begin{longtblr}[
  label = {tab:sum_gpro_gpao},
  entry = none,
  caption = {Main characteristics of the default\acs{gint}, \acs{gcohe} and \acs{gunw} products.},
]{
  width   = \linewidth,
  colspec = {X[l] X[l] X[c] X[c]},
  colsep  = 2pt,
  hlines,
  hline{1-Z} = {-}{0.05em},
}

\SetCell[r=2]{c}\textbf{Characteristic Item}
  & \SetCell[r=2]{c}\textbf{Input Type}
  & \SetCell[c=2]{c}\textbf{Mission} & \\
& & \textbf{Sentinel-1} & \textbf{ROSE-L} \\

Ground Projection  & { + COP\acs{dem} 30m \\
                      + Precise Reference Orbit}& \ding{51} & \ding{51} \\

Spacing & & 50 m x 50 m & 50 m x 50 m \\

Slicing &  & S-1 projected burst &  ROSE-L projected burst \\

Reference System &  & Cartographic: UTM/UPS &  Cartographic: UTM/UPS \\

\end{longtblr}


\subsection{\acs{cmlc}}
\label{sec:cmlc_sum}

The default \acs{cmlc} product is generated from the \acs{cslc}. No timing or atmospheric corrections are applied during \acs{cmlc} generation, as these are handled at the \acs{cslc} level. Below, only the product-specific main characteristics are reported.

\begin{longtblr}[
  label = {tab:sum_cmlc},
  entry = none,
  caption = {Main characteristics of the default \acs{cmlc} product.},
]{
  width   = \linewidth,
  colspec = {X[l] X[l] X[c] X[c]},
  colsep  = 2pt,
  hlines,
  hline{1-Z} = {-}{0.05em},
}

\SetCell[r=2]{c}\textbf{Characteristic Item}
  & \SetCell[r=2]{c}\textbf{Input Type}
  & \SetCell[c=2]{c}\textbf{Mission} & \\
& & \textbf{Sentinel-1} & \textbf{ROSE-L} \\

Phase & Stack of \acs{cslc} & \ding{51} &  \ding{51} \\

Long-term coherence & Stack of \acs{cslc} & \ding{51} &  \ding{51} \\

Short-term coherence & Stack of \acs{cslc} & \ding{51} &  \ding{51} \\

Spacing & & $\sim$50 m x 50 m & $\sim$50 m x 50 m \\

Slicing &  & S-1 burst &  ROSE-L burst \\

Reference System &  & Azimuth/Slant-range &  Azimuth/Slant-range \\

\end{longtblr}

\subsection{\acs{gmlc}}
\label{sec:gmlc_summ}

The default \acs{gmlc} is the geocoded \acs{cmlc} product, from which it inherits the main characteristics defined in \ref{sec:cmlc_sum}. Below, only the product-specific main characteristics are reported.

\begin{longtblr}[
  label = {tab:sum_gmlc},
  entry = none,
  caption = {Main characteristics of the default \acs{gmlc} product.},
]{
  width   = \linewidth,
  colspec = {X[l] X[l] X[c] X[c]},
  colsep  = 2pt,
  hlines,
  hline{1-Z} = {-}{0.05em},
}

\SetCell[r=2]{c}\textbf{Characteristic Item}
  & \SetCell[r=2]{c}\textbf{Input Type}
  & \SetCell[c=2]{c}\textbf{Mission} & \\
& & \textbf{Sentinel-1} & \textbf{ROSE-L} \\

Phase & Stack of \acs{cslc} & \ding{51} &  \ding{51} \\

Long-term coherence & Stack of \acs{cslc} & \ding{51} &  \ding{51} \\

Short-term coherence & Stack of \acs{cslc} & \ding{51} &  \ding{51} \\


Ground Projection  & { + COP\acs{dem} 30m \\
                     + Precise Reference Orbit}& \ding{51} & \ding{51} \\

Spacing & & 50 m x 50 m & 50 m x 50 m \\

Slicing &  & S-1 projected burst &  ROSE-L projected burst \\

Reference System &  & Cartographic: UTM/UPS &  Cartographic: UTM/UPS \\

\end{longtblr}

\subsection{ENRB}
\label{sec:enh_nrb_summ}

The default \acs{enrb} is the version of the \acs{nrb} product, generated by applying specific filtering procedures to the input \acs{cslc}. Below, only the product-specific main characteristics are reported.

\begin{longtblr}[
  label = {tab:sum_enrb},
  entry = none,
  caption = {Main characteristics of the default \acs{enrb} product.},
]{
  width   = \linewidth,
  colspec = {X[l] X[l] X[c] X[c]},
  colsep  = 2pt,
  hlines,
  hline{1-Z} = {-}{0.05em},
}

\SetCell[r=2]{c}\textbf{Characteristic Item}
  & \SetCell[r=2]{c}\textbf{Input Type}
  & \SetCell[c=2]{c}\textbf{Mission} & \\
& & \textbf{Sentinel-1} & \textbf{ROSE-L} \\

Amplitude convention &  & gamma0 & gamma0 \\

\acs{rtc} & { + \acs{cslc} \\
              + COP\acs{dem} 30} m & \ding{51} & \ding{51} \\

Shiroma for \acs{rtc} & & \ding{51} & \ding{51} \\

Ghost Removal & & \ding{51} & \ding{51} \\

Side-lobe suppression & & \ding{51} & \ding{51} \\

\acs{anf} & & \ding{51} & \ding{51} \\

Local Incidence Angle & & \ding{51} & \ding{51} \\

Spacing & & 10 m x 10 m & 10 m x 10 m \\

Multilooking &  & - &  - \\

Slicing &  & US-\acs{mgrs} 100 km x 100 km &  US-\acs{mgrs} 100 km x 100 km \\

Reference System &  & Cartographic: UTM/UPS &  Cartographic: UTM/UPS \\

\end{longtblr}

\subsection{EORB}
\label{sec:enh_orb_summ}

The default \acs{eorb} is the version of the \acs{orb} product, generated by applying specific filtering procedures to the input \acs{cslc}. Below, only the product-specific main characteristics are reported.

\begin{longtblr}[
  label = {tab:sum_eorb},
  entry = none,
  caption = {Main characteristics of the default \acs{eorb} product.},
]{
  width   = \linewidth,
  colspec = {X[l] X[l] X[c] X[c]},
  colsep  = 2pt,
  hlines,
  hline{1-Z} = {-}{0.05em},
}

\SetCell[r=2]{c}\textbf{Characteristic Item}
  & \SetCell[r=2]{c}\textbf{Input Type}
  & \SetCell[c=2]{c}\textbf{Mission} & \\
& & \textbf{Sentinel-1} & \textbf{ROSE-L} \\

 &  &  &  \\

 &  &  &  \\

 &  &  &  \\

\end{longtblr}

\subsection{GGM}
\label{sec:ggm_summ}

The default \acs{ggm} product is generated from the \acs{gpro} and \acs{gpao}. Below, only the product-specific main characteristics are reported.

\begin{longtblr}[
  label = {tab:sum_ggm},
  entry = none,
  caption = {Main characteristics of the default \acs{ggm} product.},
]{
  width   = \linewidth,
  colspec = {X[l] X[l] X[c] X[c]},
  colsep  = 2pt,
  hlines,
  hline{1-Z} = {-}{0.05em},
}

\SetCell[r=2]{c}\textbf{Characteristic Item}
  & \SetCell[r=2]{c}\textbf{Input Type}
  & \SetCell[c=2]{c}\textbf{Mission} & \\
& & \textbf{Sentinel-1} & \textbf{ROSE-L} \\

Estimated Deformations & Pairs of \acs{gpro} and \acs{gpao} from ascending and descending geometries & Northward, Eastward and vertical deformations &  Northward, Eastward and vertical deformations \\

Spacing & & 50 m x 50 m & 50 m x 50 m \\

Slicing &  & US-\acs{mgrs} 100 km x 100 km &  US-\acs{mgrs} 100 km x 100 km \\

Reference System &  & Cartographic: UTM/UPS &  Cartographic: UTM/UPS \\

\end{longtblr}

\subsection{Stripmap equivalent CSLC}
\label{sec:ggm_summ}

The default Stripmap-equivalent (\acs{stripeq}) \acs{cslc} image is a full-resolution SAR product that preserves the native resolution of the input acquisition. It stitches multiple bursts from wide-swath modes (e.g., TOPSAR, ScanSAR) into a single image while maintaining spatial consistency and phase continuity characteristic of Stripmap mode. The output retains all the information and main characteristics of the input \acs{cslc} (see also \ref{sec:cslc_summ}).

\begin{longtblr}[
  label = {tab:sum_strcslc},
  entry = none,
  caption = {Main characteristics of the default Stripmap-equivalent \acs{cslc} product.},
]{
  width   = \linewidth,
  colspec = {X[l] X[l] X[c] X[c]},
  colsep  = 2pt,
  hlines,
  hline{1-Z} = {-}{0.05em},
}

\SetCell[r=2]{c}\textbf{Characteristic Item}
  & \SetCell[r=2]{c}\textbf{Input Type}
  & \SetCell[c=2]{c}\textbf{Mission} & \\
& & \textbf{Sentinel-1} & \textbf{ROSE-L} \\

Radar Reference Grid & & Virtual & Virtual \\

Geometric alignment & { + COP\acs{dem} 30 m \\
                        + Precise Orbits } & \ding{51} & \ding{51} \\

Topographic phase correction & COP\acs{dem} 30 m & \ding{51} & \ding{51} \\

Atmospheric delay \& timing corrections & \acs{etad} & \ding{51} & \ding{51} \\

ESD & Geometric-corrected \acs{cslc}s (GC-CSLC) & - & \ding{51} \\

ICC & Azimuth-shift-corrected GC-\acs{cslc}s (ASGC-CSLC) & - & \ding{51} \\

Split-Spectrum & { + Range-Azimuth-shift-corrected GC-\acs{cslc}s (RASGC-CSLC) \\
                   + Reference ASGC-\acs{cslc} }& - & \ding{51} \\

\SetCell[r=2]{l}Spacing
  & \SetCell[r=2]{l}
  & \SetCell[r=2,c=2]{c}
    {Defined by the Radar Reference Grid:\\
     minimum between Sentinel-1\\
     and ROSE-L spacing} & \\
& & & \\

Multilooking &  & - &  - \\

Denoising &  & - &  - \\

Slicing &  & S-1 product tile &  ROSE-L product tile \\

Reference System &  & Azimuth/Slant-range &  Azimuth/Slant-range \\

\end{longtblr}

\end{longtblr}